\documentclass[10pt,letterpaper]{article}

\usepackage[
  letterpaper,
  top=0.7in,
  bottom=0.7in,
  left=0.8in,
  right=0.8in
]{geometry}
\usepackage[T1]{fontenc}
\usepackage{mathptmx}
\usepackage{amsmath}
\usepackage{amssymb}
\usepackage{graphicx}
\usepackage{booktabs}
\usepackage{array}
\usepackage{cite}
\usepackage{url}
\usepackage{float}

\setlength{\parindent}{0.14in}
\setlength{\parskip}{0pt}
\pagestyle{plain}
\renewcommand{\thesection}{\Roman{section}}
\renewcommand{\thesubsection}{\Alph{subsection}.}
\renewcommand{\thetable}{\Roman{table}}
\renewcommand{\refname}{REFERENCES}

\makeatletter
\renewcommand\section{\@startsection{section}{1}{\z@}%
  {1.3ex plus 1ex minus .2ex}%
  {0.8ex plus .2ex}%
  {\centering\normalfont\small\scshape}}
\renewcommand\subsection{\@startsection{subsection}{2}{\z@}%
  {1.0ex plus .8ex minus .2ex}%
  {0.6ex plus .2ex}%
  {\normalfont\itshape}}
\renewcommand{\fnum@table}{\textsc{Table \thetable}}
\makeatother

\begin{document}
\twocolumn[

\begin{center}
{\LARGE An investigation into Lightweight Hand Pose estimation for Edge HCI\par}
\end{center}

\vspace{0.6em}
\noindent\textbf{\textit{Abstract}}---This paper investigates lightweight 2D RGB hand pose estimation for edge-based human-computer interaction using a DRHand-inspired dual-branch architecture. The study focuses on three questions: whether a reduced 10-keypoint representation can retain sufficient utility for interaction-oriented use, whether dual-branch fusion improves prediction relative to the individual heatmap and coordinate branches, and whether pseudo-labeled real-image data can improve generalization beyond synthetic-only training. The implemented pipeline is evaluated through controlled experiments on RHD and a COCO-style hand dataset, together with transfer-learning runs between the two domains, comparison against a MediaPipe baseline, deployment-oriented benchmarking on Jetson Orin Nano, and qualitative analysis on selected benchmark and real-image examples. The results show that the reduced 10-keypoint model preserves most useful predictive performance while offering a modest runtime advantage, indicating that full 21-keypoint prediction is not always necessary for lightweight edge-HCI settings. Fusion improves the overall balance of results, although diagnostic analysis shows that the current heuristic remains suboptimal. Transfer learning substantially improves cross-domain performance when used for target-domain adaptation, but zero-shot generalization remains weak. The learned DRHand variants outperform the evaluated MediaPipe baseline on benchmark RHD and achieve practical embedded runtime. Overall, the study establishes a strong RGB baseline for lightweight hand-pose estimation on edge devices and identifies improved fusion, broader real-image evaluation, and detector-based input normalization as key directions for future work.

\vspace{0.6em}
\noindent\textbf{\textit{Keywords}}---hand pose estimation; human-computer interaction; edge devices; lightweight deep learning; transfer learning; embedded vision
\vspace{0.8em}
]

\section{INTRODUCTION}
\noindent\textbf{Point:} Hand pose estimation is a relevant interface mechanism for touchless HCI.
\textbf{Evidence:} Prior HCI research describes hand-based interaction as natural, expressive, and useful for pointing, manipulation, and contact-free control in settings such as AR/VR, large displays, and clinical environments~\cite{vuletic,wachs,li_vr}.
\textbf{Explanation:} Unlike simple gesture classification, landmark-based hand pose provides a structured representation that can support richer interaction.
\textbf{Link:} This makes hand-pose tracking a suitable target for practical HCI research.

\medskip
\noindent\textbf{Point:} The main technical challenge is achieving useful hand-pose performance under edge-device constraints.
\textbf{Evidence:} Real deployment must balance prediction quality with latency, throughput, and computational cost on constrained hardware rather than desktop GPUs~\cite{attentionlite,srhandnet,dualreg,stereojetson}.
\textbf{Explanation:} A highly accurate model is not sufficient for HCI if it cannot respond quickly enough in embedded settings.
\textbf{Link:} The paper therefore focuses on lightweight RGB hand-pose estimation for edge deployment.

\medskip
\noindent\textbf{Point:} Existing lightweight methods show that real-time hand tracking is feasible, but they leave open design questions about representation, fusion, and generalisation.
\textbf{Evidence:} MediaPipe Hands demonstrates efficient on-device RGB tracking, while DRHand combines heatmap and coordinate regression in a lightweight dual-branch framework~\cite{mediapipe,dualreg}.
\textbf{Explanation:} These methods establish a strong baseline, but they do not by themselves show whether fewer landmarks are sufficient, whether dual-branch fusion is necessary, or how robust such models are across different training domains.
\textbf{Link:} This motivates a targeted study of reduced keypoint prediction, branch utility, and cross-domain transfer.

\medskip
\noindent\textbf{Point:} This paper investigates those questions through a DRHand-inspired implementation and a deployment-oriented evaluation.
\textbf{Evidence:} The study compares 21-keypoint and reduced 10-keypoint prediction, evaluates fusion against the individual branches, analyses transfer between RHD and COCO-style data, compares against MediaPipe, and benchmarks runtime on Jetson-class hardware.
\textbf{Explanation:} These experiments are designed to test both prediction quality and practical deployment tradeoffs in one pipeline.
\textbf{Link:} The contribution of the paper is therefore not only an implementation of a lightweight hand-pose model, but also an empirical study of how such a model behaves under realistic edge-HCI constraints.

\section{LITERATURE REVIEW}.

\medskip
\noindent\textbf{Point:} The literature does not converge on a single architectural strategy for efficient pose estimation.  
\textbf{Evidence:} Existing approaches include CNN-based systems such as MediaPipe Hands~\cite{mediapipe}, lightweight attention-based models~\cite{attentionlite}, hybrid CNN-transformer designs such as EdgePose~\cite{edgepose}, temporal GCN-based tracking~\cite{handobjectgcn}, and encoder-decoder approaches such as EfficientPose~\cite{efficientpose}.  
\textbf{Explanation:} This suggests that efficient pose estimation remains an open design problem, with different model families making different tradeoffs between local feature extraction, global context, and structured joint reasoning.  
\textbf{Link:} Accordingly, it is reasonable to evaluate a DRHand-inspired architecture rather than assume that the design space has already been settled.

\medskip
\noindent\textbf{Point:} Edge deployment imposes a genuine speed-accuracy tradeoff in pose estimation.  
\textbf{Evidence:} Prior work notes that high-accuracy egocentric and 3D pose systems often rely on heavy backbones and desktop-class computation~\cite{everydayegohands}, whereas deployment-oriented methods such as MediaPipe Hands~\cite{mediapipe} and SPLite Hand~\cite{splitehand} pursue simplified architectures for constrained devices.  
\textbf{Explanation:} The literature therefore shows a persistent mismatch between models optimised for benchmark accuracy and models that are practical on embedded hardware.  
\textbf{Link:} This justifies treating runtime and deployment cost as core evaluation criteria rather than secondary implementation concerns.

\medskip
\noindent\textbf{Point:} Lightweight pose estimation methods typically achieve efficiency through deliberate simplification of the pipeline.  
\textbf{Evidence:} Reported strategies include compact backbones and lightweight attention~\cite{attentionlite}, sparse or reduced-complexity computation~\cite{splitehand}, simplified decoder heads~\cite{efficientpose}, and detector-to-tracker pipelines with cropped regions of interest~\cite{mediapipe}.  
\textbf{Explanation:} These studies imply that efficient performance usually comes from reducing unnecessary computation across multiple stages of the model rather than from a single isolated modification.  
\textbf{Link:} This provides literature support for testing whether simpler prediction targets or prediction strategies can preserve useful hand-pose performance.

\medskip
\noindent\textbf{Point:} The choice of training data is likely to affect hand-pose model generalisation.  
\textbf{Evidence:} Prior work uses both synthetic datasets such as RHD and real datasets for hand and hand-object tracking, while augmentation is often introduced to improve robustness under more realistic conditions~\cite{handobjectgcn,lightweightgcnhand,handaugment}.  
\textbf{Explanation:} These data sources differ in appearance, pose distribution, and annotation conditions, so a model trained on one domain may not generalise reliably to another. This is especially relevant when a study relies heavily on synthetic supervision but aims to operate on real camera input.  
\textbf{Link:} Therefore, dataset realism and cross-domain robustness remain important considerations in lightweight hand-pose estimation.

\medskip
\noindent\textbf{Point:} The literature leaves open whether full 21-keypoint hand representations are always necessary for lightweight edge-HCI use.  
\textbf{Evidence:} Existing lightweight hand-pose systems such as MediaPipe Hands~\cite{mediapipe}, Attention!~\cite{attentionlite}, and SPLite Hand~\cite{splitehand} focus mainly on reducing model complexity while still predicting a full hand structure.  
\textbf{Explanation:} This suggests that prior work has concentrated more on compressing the model than on questioning whether the full landmark target is required for practical HCI interaction.  
\textbf{Link:} This directly motivates the present study's comparison between reduced 10-keypoint and standard 21-keypoint prediction.

\medskip
Taken together, these gaps motivate the three main questions examined in this study: whether a reduced 10-keypoint representation is sufficient for lightweight edge HCI, whether the added complexity of dual-branch fusion is justified, and whether pseudo-labeled real-image data can improve generalisation beyond synthetic-only training.


\section{METHOD}

\subsection{Network Structure}

\noindent\textbf{Point:} The proposed model uses a shared lightweight backbone followed by two prediction branches.  
\textbf{Explanation:} The input is a resized RGB hand image, and most computation is shared before the network splits into heatmap and coordinate regression paths. This keeps the model compact while still allowing two complementary prediction strategies. The same structure is used for both the standard 21-keypoint model and the reduced 10-keypoint variant, with only the output dimensionality changing.

Formally, the input to the model is an RGB image resized to $256 \times 256$, denoted by $x \in \mathbb{R}^{3 \times 256 \times 256}$. The shared backbone then produces a compact feature representation,
\[
F = B(x).
\]

\subsubsection{The Backbone Network}

\noindent\textbf{Point:} The backbone is designed to extract useful hand features at low computational cost.  
\textbf{Explanation:} In the implementation, the backbone is built from depthwise-separable convolutions, downsampling blocks, and skip-connected upsampling stages. This shared feature extractor is reused by both prediction branches, which makes it the efficiency-oriented core of the model.

\subsubsection{The Heatmap Regression Branch}

\noindent\textbf{Point:} The heatmap branch predicts one spatial confidence map per keypoint.  
\textbf{Explanation:} This branch takes the shared backbone features and produces a stack of $64 \times 64$ heatmaps, allowing the model to represent likely landmark locations spatially. The branch is supervised using Gaussian-target mean-squared error, which matches the heatmap output format.

From the shared representation, the heatmap branch predicts one heatmap per keypoint,
\[
\hat{H} = R_h(F), \qquad \hat{H} \in \mathbb{R}^{K \times 64 \times 64},
\]
where $K$ is the number of predicted keypoints.

For the heatmap branch, each ground-truth keypoint $p_i = (x_i, y_i)$ is converted into a Gaussian target heatmap,
\[
G_i(u,v) = \exp\left(
-\frac{(u-x_i)^2 + (v-y_i)^2}{2\sigma^2}
\right),
\]
and the branch is supervised with mean-squared error,
\[
\mathcal{L}_{hm}
=
\frac{1}{KHW}
\sum_{i=1}^{K}
\sum_{u,v}
\left(\hat{H}_i(u,v) - G_i(u,v)\right)^2.
\]

\subsubsection{The Coordinate Regression Branch}

\noindent\textbf{Point:} The coordinate branch predicts landmark positions directly from intermediate heatmap features.  
\textbf{Explanation:} Instead of producing spatial maps, this branch regresses normalized $(x,y)$ coordinates for each keypoint. It operates on the intermediate features from the heatmap path, which keeps the branches related rather than fully independent. The branch is supervised using Wing loss, which is suited to landmark regression.

The coordinate branch predicts normalized landmark coordinates,
\[
\hat{P}^{coord} = R_c(F_{64}), \qquad \hat{P}^{coord} \in \mathbb{R}^{K \times 2},
\]
where $F_{64}$ denotes the intermediate feature map passed from the heatmap path into the coordinate regressor.

The coordinate branch is supervised directly using Wing loss,
\[
\mathcal{L}_{coord}
=
\frac{1}{2K}
\sum_{i=1}^{K}
\sum_{j \in \{x,y\}}
\phi\!\left(\hat{p}^{coord}_{i,j} - p_{i,j}\right),
\]
where
\[
\phi(\delta)=
\begin{cases}
w \ln\left(1 + \frac{|\delta|}{\epsilon}\right), & |\delta| < w, \\
|\delta| - C, & \text{otherwise,}
\end{cases}
\qquad
C = w - w \ln\left(1 + \frac{w}{\epsilon}\right).
\]

\subsection{Post Processing}

\noindent\textbf{Point:} The final pose can be produced by either branch alone or by a simple fusion rule.  
\textbf{Explanation:} The implementation supports heatmap-only, coordinate-only, and fusion inference modes. In fusion mode, the heatmaps are decoded by argmax and compared with the coordinate predictions on a per-joint basis. The final decision is made using a threshold derived from the median predicted bone length, so the post-processing remains simple while still combining the strengths of both branches.

At inference time, the heatmap prediction for keypoint $i$ is decoded by argmax. If $(u_i^*, v_i^*)$ is the index of the peak response in $\hat{H}_i$, then the decoded heatmap coordinate is
\[
\hat{p}^{hm}_i =
\left(
\frac{u_i^*}{W-1},
\frac{v_i^*}{H-1}
\right).
\]
For fusion, the disagreement between the two predictions for joint $i$ is
\[
d_i = \left\| \hat{p}^{hm}_i - \hat{p}^{coord}_i \right\|_2.
\]
A sample-specific threshold $\alpha$ is computed as the median predicted bone length,
\[
\alpha = \mathrm{median}_{(a,b)\in \mathcal{E}}
\left\| \hat{p}^{coord}_a - \hat{p}^{coord}_b \right\|_2,
\]
where $\mathcal{E}$ is the set of hand-skeleton edges. The fused output is then defined as
\[
\hat{p}^{fused}_i =
\begin{cases}
\hat{p}^{hm}_i, & d_i < \alpha, \\
\hat{p}^{coord}_i, & \text{otherwise.}
\end{cases}
\]

\section{EXPERIMENTS}

\subsection{Experimental Setup}

\noindent\textbf{Point:} The experiments use two data sources: RHD and a COCO-style hand dataset.  
\textbf{Explanation:} RHD provides the main synthetic benchmark used throughout the study, while the COCO-style dataset is included as a real-image source with pseudo-labeled hand annotations. This allows the experiments to compare performance within the synthetic domain and under transfer to a more realistic image distribution.

\medskip
\noindent\textbf{Point:} Two landmark layouts are evaluated.  
\textbf{Explanation:} The standard configuration predicts 21 hand keypoints, while the reduced configuration predicts 10 tip/base keypoints. This allows the study to test whether a smaller landmark set can preserve sufficient pose information while reducing computational cost.

\medskip
\noindent\textbf{Point:} The study includes both direct training and transfer-learning experiments.  
\textbf{Explanation:} The main transfer sequences are \texttt{R\_ONLY}, \texttt{C\_ONLY}, \texttt{C\_TO\_R}, and \texttt{R\_TO\_C}. These runs are used to examine how well models trained on one domain perform on the other and how much target-domain fine-tuning improves adaptation.

\medskip
\noindent\textbf{Point:} Internal and external comparisons are both included in the setup.  
\textbf{Explanation:} Internal comparisons use the same trained checkpoints to evaluate fusion, heatmap-only, and coordinate-only prediction modes. In addition, MediaPipe is evaluated as an external baseline on RHD, and embedded inference is benchmarked separately on Jetson hardware.

\medskip
\noindent\textbf{Point:} Reported learned-model results are based on repeated runs.  
\textbf{Explanation:} Unless otherwise stated, results for the learned DRHand variants are summarized over five random seeds and reported as mean $\pm$ standard deviation. This reduces the likelihood that conclusions depend on a single training run.

\subsection{Implementation Details}

\noindent\textbf{Point:} The proposed system is implemented in PyTorch.  
\textbf{Explanation:} This implementation follows the DRHand-inspired dual-branch design described earlier, but is built as a configurable PyTorch training and evaluation framework rather than as a direct TensorFlow reproduction of the original paper.

\medskip
\noindent\textbf{Point:} All model inputs are normalized RGB images resized to a fixed spatial resolution.  
\textbf{Explanation:} In the present implementation, images are resized to $256 \times 256$ before being passed to the network. This fixed input size is used consistently across training and evaluation to keep the model configuration stable.

\medskip
\noindent\textbf{Point:} Training follows a two-stage DRHand-style procedure.  
\textbf{Explanation:} In stage 1, the backbone and heatmap branch are optimized first. In stage 2, the coordinate branch is trained together with optional fine-tuning of the rest of the model. This preserves the staged training logic of the original DRHand design while allowing the branches to be optimized sequentially.

\medskip
\noindent\textbf{Point:} Transfer-learning runs are implemented through checkpoint initialization and target-domain fine-tuning.  
\textbf{Explanation:} For \texttt{C\_TO\_R} and \texttt{R\_TO\_C}, training begins from a source-domain checkpoint and continues on the target dataset. This makes transfer learning part of the same implementation framework rather than a separate training pipeline.

\medskip
\noindent\textbf{Point:} Training and checkpoint-based evaluation were performed on high-performance GPU hardware.  
\textbf{Explanation:} The learned models were trained and evaluated on an NVIDIA Tesla A100 80GB GPU. This provided the main development environment for the experiments reported in the paper.

\medskip
\noindent\textbf{Point:} Practical implementation choices were made where needed to accommodate resource constraints and incomplete source details.  
\textbf{Explanation:} Some training parameters were not fully specified in the original DRHand paper, so explicit reconstruction choices were required in this implementation. In addition, gradient accumulation was used where necessary to manage memory constraints during training.

\medskip
\noindent\textbf{Point:} Embedded deployment performance was measured separately from training.  
\textbf{Explanation:} While model development was carried out on the A100 system, runtime benchmarking for deployment-oriented evaluation was performed separately on Jetson Orin Nano hardware.

\subsection{Evaluation Metrics}

\noindent\textbf{Point:} Three main pose-estimation metrics are used to evaluate prediction quality.  
\textbf{Explanation:} The learned models are evaluated using sum of squared error (SSE), end-point error (EPE), and probability of correct keypoint (PCK). These metrics are computed on normalized coordinates and only over visible evaluated keypoints. Let $v_{si} \in \{0,1\}$ denote the visibility of keypoint $i$ in sample $s$.

The normalized SSE is defined as
\[
\mathrm{SSE} =
\frac{1}{N}
\sum_{s=1}^{N}
\sum_{i=1}^{K}
v_{si}\left\|
\hat{p}_{si} - p_{si}
\right\|_2^2,
\]
where $N$ is the number of samples, $K$ is the number of evaluated keypoints, $p_{si}$ is the ground-truth coordinate for keypoint $i$ in sample $s$, and $\hat{p}_{si}$ is the predicted coordinate.

\medskip
\noindent\textbf{Point:} End-point error is measured in a root-relative normalized form.  
\textbf{Explanation:} EPE is computed after subtracting the root joint from both the prediction and ground truth, which reduces the effect of global translation and measures the relative landmark displacement within the hand pose. Let $r$ denote the root keypoint used for the current evaluation layout. In the full layout, keypoint $0$ is used when available. In reduced or shared-10 evaluation, keypoint $1$ is used when keypoint $0$ is absent.

The normalized root-relative EPE is defined as
\[
\mathrm{EPE} =
\frac{
\sum_{s=1}^{N}
\sum_{i=1}^{K}
v_{si}
\left\|
(\hat{p}_{si} - \hat{p}_{sr}) - (p_{si} - p_{sr})
\right\|_2
}{
\sum_{s=1}^{N}
\sum_{i=1}^{K}
v_{si}
},
\]
where EPE is reported only when a supported root keypoint is present in the evaluated layout.

\medskip
\noindent\textbf{Point:} PCK measures the proportion of visible keypoints localized within a fixed error threshold.  
\textbf{Explanation:} In this study, a keypoint is counted as correct when its normalized distance from the ground truth is at most $0.2$. Higher PCK therefore indicates a larger fraction of acceptable landmark predictions.

\[
\mathrm{PCK}_{\sigma} =
\frac{
\sum_{s=1}^{N}
\sum_{i=1}^{K}
\;v_{si}\,\mathbf{1}
\left(
\left\|
\hat{p}_{si} - p_{si}
\right\|_2 \le \sigma
\right)
}{
\sum_{s=1}^{N}
\sum_{i=1}^{K}
v_{si}
},
\qquad \sigma = 0.2,
\]
where $\mathbf{1}(\cdot)$ is the indicator function.

\medskip
\noindent\textbf{Point:} Runtime is reported separately from pose accuracy.  
\textbf{Explanation:} Deployment-oriented evaluation includes prediction time in milliseconds per image and throughput in images per second. For the Jetson measurements, both forward-only and end-to-end timings are reported to distinguish model computation from surrounding pipeline overhead.

\medskip
\noindent\textbf{Point:} Additional metrics are used for baseline and fusion-specific analysis.  
\textbf{Explanation:} For the MediaPipe baseline, detection success rate is reported because the pipeline may fail before landmark prediction. For fusion analysis, the study also reports branch-selection rates, disagreement between branch outputs, the mean fusion threshold $\alpha$, and the oracle-match rate, which measures how often the heuristic selected the lower-error branch.

\section{RESULTS AND ANALYSIS}

\subsection{Experimental Results}

\subsubsection{Reduced Keypoint Performance}

\noindent\textbf{Point:} The reduced 10-keypoint model retains broadly comparable performance to the 21-keypoint baseline on the shared subset.  
\textbf{Evidence:} In Table~\ref{tab:main_accuracy}, the shared-10 evaluation of \texttt{B0} gives SSE 0.1511, EPE 0.0764, and PCK 0.9101, while the native-10 \texttt{B1} model gives SSE 0.1574, EPE 0.0706, and PCK 0.9099.  
\textbf{Explanation:} These values are very close, indicating that reducing the landmark set does not substantially weaken performance on the subset of keypoints most relevant to the reduced representation.  
\textbf{Link:} This supports the claim that a smaller landmark set can remain viable for lightweight edge-HCI use.

\begin{table*}[t]
\centering
\footnotesize
\caption{Main quantitative comparison between the 21-keypoint baseline and the reduced 10-keypoint model on RHD. Lower SSE and EPE are better; higher PCK@0.2 is better. Values are mean $\pm$ std over five seeds.}
\label{tab:main_accuracy}
\begin{tabular}{llcccc}
\toprule
Model & Eval set & Train joints & SSE $\downarrow$ & EPE $\downarrow$ & PCK@0.2 $\uparrow$ \\
\midrule
B0 baseline & native-21 & 21 & 0.3098 $\pm$ 0.0236 & 0.0779 $\pm$ 0.0054 & 0.9112 $\pm$ 0.0071 \\
B0 baseline & shared-10 & 21 & 0.1511 $\pm$ 0.0120 & 0.0764 $\pm$ 0.0039 & 0.9101 $\pm$ 0.0071 \\
B1 reduced & native-10 & 10 & 0.1574 $\pm$ 0.0226 & 0.0706 $\pm$ 0.0054 & 0.9099 $\pm$ 0.0046 \\
\bottomrule
\end{tabular}
\end{table*}

\subsubsection{Runtime and Deployment Cost}

\medskip
\noindent\textbf{Point:} The runtime advantage of the 10-keypoint model appears mainly in fusion mode.  
\textbf{Evidence:} In Table~\ref{tab:latency}, the forward-time gap between \texttt{B0} and \texttt{B1} is much larger in fusion mode (27.43 vs 24.59 ms/image) than in heatmap mode (24.28 vs 24.00 ms/image) or coordinate mode (24.08 vs 23.91 ms/image).  
\textbf{Explanation:} Since the single-branch modes differ only slightly, the main efficiency gain of the reduced model appears to come from lower fusion-stage overhead rather than from a large reduction in shared network cost.  
\textbf{Link:} This suggests that the practical speed benefit of the 10-keypoint model is tied mainly to the post-processing stage.

\begin{table*}[t]
\centering
\footnotesize
\caption{Edge-device latency summary for the main models. Forward time measures the prediction stage only, whereas end-to-end time includes image read, preprocessing, host-to-device transfer, prediction, and JSON write. Values are mean $\pm$ std over five seeds.}
\label{tab:latency}
\begin{tabular}{llccc}
\toprule
Model & Mode & Forward ms/img $\downarrow$ & End-to-end ms/img $\downarrow$ & Images/s $\uparrow$ \\
\midrule
B0 & fusion & 27.43 $\pm$ 0.45 & 36.81 $\pm$ 1.33 & 27.20 $\pm$ 0.97 \\
B0 & heatmap & 24.28 $\pm$ 0.35 & 33.19 $\pm$ 0.56 & 30.13 $\pm$ 0.50 \\
B0 & coord & 24.08 $\pm$ 0.44 & 33.17 $\pm$ 0.86 & 30.16 $\pm$ 0.77 \\
B1 & fusion & 24.59 $\pm$ 0.61 & 33.31 $\pm$ 0.71 & 30.03 $\pm$ 0.64 \\
B1 & heatmap & 24.00 $\pm$ 0.34 & 32.64 $\pm$ 0.46 & 30.64 $\pm$ 0.43 \\
B1 & coord & 23.91 $\pm$ 0.29 & 32.50 $\pm$ 0.34 & 30.77 $\pm$ 0.32 \\
\bottomrule
\end{tabular}
\end{table*}

\begin{table*}[t]
\centering
\footnotesize
\caption{Detailed latency breakdown for the main fused models. These metrics are available from the benchmark pipeline and help localize where runtime is spent. Values are mean $\pm$ std over five seeds.}
\label{tab:latency_breakdown}
\begin{tabular}{lcccccc}
\toprule
Model & Read ms & Prep ms & H2D ms & Forward ms & Write ms & E2E ms \\
\midrule
B0 fusion & 4.89 $\pm$ 0.64 & 2.80 $\pm$ 0.13 & 0.32 $\pm$ 0.02 & 27.43 $\pm$ 0.45 & 0.67 $\pm$ 0.09 & 36.81 $\pm$ 1.33 \\
B1 fusion & 4.49 $\pm$ 0.08 & 2.74 $\pm$ 0.05 & 0.31 $\pm$ 0.01 & 24.59 $\pm$ 0.61 & 0.52 $\pm$ 0.02 & 33.31 $\pm$ 0.71 \\
\bottomrule
\end{tabular}
\end{table*}

\subsubsection{Branch Behaviour}

\medskip
\noindent\textbf{Point:} Fusion provides the strongest overall balance of branch behaviour.  
\textbf{Evidence:} In Table~\ref{tab:branch_ablation}, fusion gives the best or near-best overall combination of SSE and EPE for both \texttt{B0} and \texttt{B1}, whereas the heatmap-only and coordinate-only modes are less consistent across metrics.  
\textbf{Explanation:} This indicates that the two branches contribute complementary strengths, so the dual-regression design offers more than simply duplicating two alternative outputs.  
\textbf{Link:} Accordingly, the branch-ablation results support the use of fusion as the main operating mode of the system.

\begin{table*}[t]
\centering
\footnotesize
\caption{Branch ablation for the dual-regression design. Fusion is compared against the heatmap-only and coordinate-only outputs using the same trained checkpoints. Values are mean $\pm$ std over five seeds.}
\label{tab:branch_ablation}
\begin{tabular}{llcccc}
\toprule
Model & Eval set & Mode & SSE $\downarrow$ & EPE $\downarrow$ & PCK@0.2 $\uparrow$ \\
\midrule
B0 & native-21 & fusion & 0.3098 $\pm$ 0.0236 & 0.0779 $\pm$ 0.0054 & 0.9112 $\pm$ 0.0071 \\
B0 & native-21 & heatmap & 0.3139 $\pm$ 0.0684 & 0.0989 $\pm$ 0.0335 & 0.9374 $\pm$ 0.0066 \\
B0 & native-21 & coord & 0.3134 $\pm$ 0.0229 & 0.0788 $\pm$ 0.0062 & 0.9113 $\pm$ 0.0071 \\
B1 & native-10 & fusion & 0.1574 $\pm$ 0.0226 & 0.0706 $\pm$ 0.0054 & 0.9099 $\pm$ 0.0046 \\
B1 & native-10 & heatmap & 0.2213 $\pm$ 0.0880 & 0.0991 $\pm$ 0.0288 & 0.9022 $\pm$ 0.0393 \\
B1 & native-10 & coord & 0.1599 $\pm$ 0.0218 & 0.0715 $\pm$ 0.0070 & 0.9102 $\pm$ 0.0048 \\
\bottomrule
\end{tabular}
\end{table*}

\subsubsection{Transfer and Real-Image Adaptation}

\medskip
\noindent\textbf{Point:} The experiments show that synthetic-only training does not generalize reliably to more realistic image domains, but pseudo-labeled real-image supervision provides a useful adaptation bridge.  
\textbf{Evidence:} In Table~\ref{tab:transfer_learning}, the RHD-trained model performs strongly on RHD but poorly on the COCO-style dataset, while the COCO-trained model shows the opposite pattern. After target-domain adaptation, \texttt{R\_TO\_C} and \texttt{C\_TO\_R} both improve substantially over their cross-domain starting points. This interpretation is further supported by the pilot real-image evaluation in Table~\ref{tab:hand_frame_small_pilot}, where the COCO-oriented models, especially \texttt{C\_ONLY} and \texttt{R\_TO\_C}, perform more strongly than the RHD-only baselines.  
\textbf{Explanation:} These results suggest that the main difficulty is not simply the absence of pretraining, but the domain gap between synthetic training data and more realistic image conditions. Pseudo-labeled real-image data does not provide perfect supervision, but it appears to move the model closer to the real-image domain than synthetic-only training can achieve on its own.  
\textbf{Link:} Therefore, the transfer results are best interpreted as evidence that pseudo-labeled real-image supervision can act as a practical adaptation bridge beyond synthetic-only training.

\begin{table*}[t]
\centering
\footnotesize
\caption{Transfer-learning results for the four training sequences evaluated on both RHD and the COCO-style hand dataset. Values are mean $\pm$ std over five seeds, and fusion-mode evaluation is shown because it is the main operating mode in the study.}
\label{tab:transfer_learning}
\resizebox{\textwidth}{!}{%
\begin{tabular}{llcccccc}
\toprule
Training sequence & Exp. & RHD SSE & RHD EPE & RHD PCK & COCO SSE & COCO EPE & COCO PCK \\
\midrule
RHD scratch & \texttt{R\_ONLY} & 0.3091 $\pm$ 0.0269 & 0.0764 $\pm$ 0.0030 & 0.9143 $\pm$ 0.0097 & 3.7488 $\pm$ 0.7985 & 0.3365 $\pm$ 0.0075 & 0.3719 $\pm$ 0.0434 \\
COCO scratch & \texttt{C\_ONLY} & 1.4194 $\pm$ 0.2619 & 0.1609 $\pm$ 0.0207 & 0.3397 $\pm$ 0.0718 & 0.0764 $\pm$ 0.0016 & 0.0612 $\pm$ 0.0049 & 0.9834 $\pm$ 0.0007 \\
COCO $\rightarrow$ RHD & \texttt{C\_TO\_R} & 0.4504 $\pm$ 0.0162 & 0.1020 $\pm$ 0.0039 & 0.8190 $\pm$ 0.0114 & 1.6426 $\pm$ 0.2293 & 0.2678 $\pm$ 0.0114 & 0.5320 $\pm$ 0.0164 \\
RHD $\rightarrow$ COCO & \texttt{R\_TO\_C} & 0.5699 $\pm$ 0.0592 & 0.1276 $\pm$ 0.0149 & 0.7690 $\pm$ 0.0221 & 0.0980 $\pm$ 0.0064 & 0.0840 $\pm$ 0.0018 & 0.9837 $\pm$ 0.0007 \\
\bottomrule
\end{tabular}
}
\end{table*}

\medskip
\noindent\textbf{Point:} A small pilot real-image evaluation was also conducted on a manually annotated image set.  
\textbf{Explanation:} In addition to the main benchmark results on RHD and the COCO-style dataset, the models were evaluated on the \texttt{hand\_frame\_small} set, which contains eight manually corrected real images. These results are included as exploratory evidence of real-image behaviour rather than as a standalone benchmark because of the small sample size.

\begin{table*}[t]
\centering
\footnotesize
\caption{Pilot real-image evaluation on the manually annotated \texttt{hand\_frame\_small} set ($n=8$ images). Learned-model rows are mean $\pm$ std over five seeds; MediaPipe is a single run. All learned-model results use fusion-mode evaluation.}
\label{tab:hand_frame_small_pilot}
\begin{tabular}{llcccc}
\toprule
Method & Eval protocol & SSE $\downarrow$ & EPE $\downarrow$ & PCK@0.2 $\uparrow$ & Detect. succ. \\
\midrule
B0 & shared-10 & 0.3345 $\pm$ 0.0221 & 0.2815 $\pm$ 0.0129 & 0.7075 $\pm$ 0.0143 & n/a \\
B1 & native-10 & 0.4350 $\pm$ 0.1153 & 0.2752 $\pm$ 0.0307 & 0.6475 $\pm$ 0.0555 & n/a \\
R\_ONLY & shared-10 & 0.3595 $\pm$ 0.0475 & 0.2706 $\pm$ 0.0214 & 0.7050 $\pm$ 0.0338 & n/a \\
C\_ONLY & shared-10 & 0.1215 $\pm$ 0.0434 & 0.1155 $\pm$ 0.0286 & 0.9200 $\pm$ 0.0647 & n/a \\
C\_TO\_R & shared-10 & 0.3644 $\pm$ 0.0360 & 0.2481 $\pm$ 0.0090 & 0.7075 $\pm$ 0.0301 & n/a \\
R\_TO\_C & shared-10 & 0.1209 $\pm$ 0.0236 & 0.1043 $\pm$ 0.0151 & 0.9425 $\pm$ 0.0301 & n/a \\
MediaPipe & native-21 & 0.0012 & 0.0018 & 1.0000 & 1.0000 \\
MediaPipe & shared-10 & 0.0005 & 0.0017 & 1.0000 & 1.0000 \\
\bottomrule
\end{tabular}
\end{table*}

\subsubsection{External Baseline Comparison}

\medskip
\noindent\textbf{Point:} The learned DRHand variants outperform the off-the-shelf MediaPipe baseline on RHD.  
\textbf{Evidence:} In Table~\ref{tab:mediapipe_baseline}, MediaPipe reaches SSE 1.7049 and PCK 0.7707 on native-21 evaluation, whereas the learned \texttt{B0} fusion model reaches SSE 0.3098 and PCK 0.9112. The same pattern holds on the shared-10 comparison.  
\textbf{Explanation:} This indicates that task-specific training on the target benchmark produces substantially stronger RHD performance than the generic detector-landmarker baseline.  
\textbf{Link:} Therefore, MediaPipe remains a useful reference point, but not a competitive benchmark in this evaluation setting.

\begin{table*}[t]
\centering
\footnotesize
\caption{Comparison between the MediaPipe Hand Landmarker baseline and the learned DRHand variants on RHD. The MediaPipe rows are single evaluation runs using \texttt{--hand auto}, whereas the DRHand rows are mean $\pm$ std over five seeds. Detection success is shown only for the MediaPipe runs because the learned-model evaluations do not use the same detection stage.}
\label{tab:mediapipe_baseline}
\begin{tabular}{llcccc}
\toprule
Method & Eval set & SSE & EPE & PCK@0.2 & Detect. succ. \\
\midrule
MediaPipe & native-21 & 1.7049 & 0.0833 & 0.7707 & 0.8933 \\
MediaPipe & shared-10 & 0.8599 & 0.1378 & 0.7635 & 0.8933 \\
\texttt{B0} fusion & native-21 & 0.3098 $\pm$ 0.0236 & 0.0779 $\pm$ 0.0054 & 0.9112 $\pm$ 0.0071 & n/a \\
\texttt{B0} fusion & shared-10 & 0.1511 $\pm$ 0.0120 & 0.0764 $\pm$ 0.0039 & 0.9101 $\pm$ 0.0071 & n/a \\
\texttt{B1} fusion & native-10 & 0.1574 $\pm$ 0.0226 & 0.0706 $\pm$ 0.0054 & 0.9099 $\pm$ 0.0046 & n/a \\
\bottomrule
\end{tabular}
\end{table*}

\subsubsection{Fusion Diagnostics}

\medskip
\noindent\textbf{Point:} The current fusion heuristic is useful but clearly imperfect.  
\textbf{Evidence:} In Table~\ref{tab:fusion_diagnostics}, the oracle-match rate is only about 0.48 for \texttt{B0} and 0.53 for \texttt{B1}, even though fusion improves aggregate accuracy in the branch-ablation study.  
\textbf{Explanation:} This means that the heuristic often fails to choose the lower-error branch on a per-joint basis, but still helps overall because some of its correct decisions occur on more influential prediction cases.  
\textbf{Link:} Therefore, the current fusion rule is beneficial, but it also remains one of the clearest opportunities for future improvement.

\begin{table*}[t]
\centering
\footnotesize
\caption{Fusion diagnostics for the dual-regression post-processing rule. The oracle-match rate measures how often the final fusion choice matches the lower-error branch for visible joints. Values are mean $\pm$ std over five seeds. The table includes all currently available diagnostic exports, including a shared-10 row for \texttt{B0}.}
\label{tab:fusion_diagnostics}
\begin{tabular}{llccccc}
\toprule
Model & Eval set & HM sel. & Coord sel. & $\alpha$ mean & Disagr. mean & Oracle match \\
\midrule
B0 & native-21 & 0.426 $\pm$ 0.025 & 0.574 $\pm$ 0.025 & 0.065 $\pm$ 0.007 & 0.163 $\pm$ 0.022 & 0.481 $\pm$ 0.012 \\
B0 & shared-10 & 0.407 $\pm$ 0.024 & 0.593 $\pm$ 0.024 & 0.065 $\pm$ 0.007 & 0.162 $\pm$ 0.018 & 0.474 $\pm$ 0.011 \\
B1 & native-10 & 0.501 $\pm$ 0.053 & 0.499 $\pm$ 0.053 & 0.084 $\pm$ 0.003 & 0.171 $\pm$ 0.018 & 0.529 $\pm$ 0.037 \\
\bottomrule
\end{tabular}
\end{table*}

\subsubsection{Hyperparameter Sensitivity}

\medskip
\noindent\textbf{Point:} The baseline hyperparameter setting is a compromise rather than an obviously dominant configuration.  
\textbf{Evidence:} In Table~\ref{tab:hyper_ablation}, several alternative settings improve one metric while worsening another; for example, \texttt{E1} improves EPE and PCK relative to \texttt{B0}, whereas \texttt{W2} gives the lowest SSE, but no single configuration is uniformly best across all reported metrics.  
\textbf{Explanation:} This indicates that the baseline configuration is not trivially optimal, but instead reflects a practical tradeoff between competing accuracy criteria and runtime behaviour.  
\textbf{Link:} As a result, the hyperparameter study suggests that the selected baseline is a reasonable compromise rather than a clearly dominated operating point.

\begin{table}[H]
\centering
\footnotesize
\caption{Hyperparameter ablation relative to the baseline \texttt{B0} configuration, using native-21 fused accuracy on RHD and fused end-to-end benchmark latency. Values are mean $\pm$ std over five seeds.}
\label{tab:hyper_ablation}
\begin{tabular}{llcccc}
\toprule
Exp & Change vs. B0 & SSE $\downarrow$ & EPE $\downarrow$ & PCK@0.2 $\uparrow$ & E2E ms $\downarrow$ \\
\midrule
B0 & baseline & 0.3098 $\pm$ 0.0236 & 0.0779 $\pm$ 0.0054 & 0.9112 $\pm$ 0.0071 & 36.81 $\pm$ 1.33 \\
L1 & $\lambda_{hm}=1.5$, $\lambda_{coord}=0.5$ & 0.3583 $\pm$ 0.0534 & 0.0725 $\pm$ 0.0096 & 0.9081 $\pm$ 0.0092 & 35.77 $\pm$ 0.45 \\
L2 & $\lambda_{hm}=0.5$, $\lambda_{coord}=1.5$ & 0.3073 $\pm$ 0.0260 & 0.0812 $\pm$ 0.0022 & 0.9115 $\pm$ 0.0095 & 35.75 $\pm$ 0.49 \\
S1 & $\sigma=1.5$ & 0.3252 $\pm$ 0.0381 & 0.0707 $\pm$ 0.0055 & 0.9116 $\pm$ 0.0086 & 35.72 $\pm$ 0.25 \\
S2 & $\sigma=3.0$ & 0.3137 $\pm$ 0.0488 & 0.0833 $\pm$ 0.0058 & 0.9177 $\pm$ 0.0081 & 36.48 $\pm$ 0.53 \\
W1 & $w=5.0$ & 0.3213 $\pm$ 0.0384 & 0.0737 $\pm$ 0.0058 & 0.9150 $\pm$ 0.0099 & 36.21 $\pm$ 0.39 \\
W2 & $w=15.0$ & 0.2979 $\pm$ 0.0441 & 0.0750 $\pm$ 0.0031 & 0.9140 $\pm$ 0.0124 & 35.91 $\pm$ 0.43 \\
E1 & $\epsilon=1.0$ & 0.2994 $\pm$ 0.0591 & 0.0695 $\pm$ 0.0011 & 0.9175 $\pm$ 0.0113 & 35.84 $\pm$ 0.23 \\
E2 & $\epsilon=4.0$ & 0.3106 $\pm$ 0.0106 & 0.0873 $\pm$ 0.0056 & 0.9110 $\pm$ 0.0067 & 36.13 $\pm$ 0.52 \\
\bottomrule
\end{tabular}
\end{table}

\subsection{Qualitative Results}

\medskip
\noindent\textbf{Point:} The real-image examples favour the COCO-oriented models.  
\textbf{Evidence:} In the two \texttt{hand\_frame\_small} figures, \texttt{C\_ONLY} and especially \texttt{R\_TO\_C} place the red landmarks closer to the green annotations than \texttt{B0} and \texttt{B1}, particularly around the inner finger cluster.  
\textbf{Explanation:} This suggests that the real-image examples are visually closer to the COCO-style domain than to the synthetic RHD domain.  
\textbf{Link:} The qualitative results therefore support the transfer findings on real-image adaptation.

\medskip
\noindent\textbf{Point:} The COCO-only model generalizes poorly to RHD.  
\textbf{Evidence:} In the RHD examples \texttt{00143} and \texttt{00354}, \texttt{C\_ONLY} produces landmark groups that are visibly detached from the hand.  
\textbf{Explanation:} This is a clear visual failure under domain shift rather than a minor local error.  
\textbf{Link:} The selected figures therefore reinforce the conclusion that zero-shot COCO-to-RHD transfer is weak.

\medskip
\noindent\textbf{Point:} Target-domain adaptation improves the transferred COCO model on RHD.  
\textbf{Evidence:} In \texttt{00354}, \texttt{C\_TO\_R} is visibly better aligned to the hand than \texttt{C\_ONLY}, even though it is not the strongest prediction in the figure.  
\textbf{Explanation:} Fine-tuning reduces the most obvious cross-domain failures, but does not eliminate them completely.  
\textbf{Link:} This supports the interpretation that transfer is useful mainly as initialization for later target-domain adaptation.

\medskip
\noindent\textbf{Point:} \texttt{R\_TO\_C} is the most consistently strong learned model across the selected figures.  
\textbf{Evidence:} It is among the closest models to ground truth on the two real-image examples and remains stable on the two RHD examples, unlike \texttt{C\_ONLY}.  
\textbf{Explanation:} This indicates that RHD pretraining followed by COCO-style adaptation transfers more consistently across the displayed cases than COCO-only training.  
\textbf{Link:} The selected figures therefore visually support the transfer asymmetry seen in the quantitative results.

\medskip
\noindent\textbf{Point:} MediaPipe is visually strong, but the real-image comparison is not fully independent.  
\textbf{Evidence:} MediaPipe aligns closely with the annotations in the two \texttt{hand\_frame\_small} figures, but those annotations were initialized from MediaPipe predictions before manual correction.  
\textbf{Explanation:} This gives MediaPipe a structural advantage on that small pilot set.  
\textbf{Link:} The real-image MediaPipe comparisons should therefore be treated as illustrative rather than definitive.

\section*{QUALITATIVE FIGURES}

\begin{figure*}[t]
\centering
\includegraphics[width=\textwidth]{../qual_figures/hand_frame_small_best/WIN_20260405_12_56_28_Pro_qualitative.png}
\caption{Qualitative comparison on hand\_frame\_small example 1.}
\end{figure*}

\begin{figure*}[t]
\centering
\includegraphics[width=\textwidth]{../qual_figures/hand_frame_small_best/WIN_20260405_12_56_57_Pro_qualitative.png}
\caption{Qualitative comparison on hand\_frame\_small example 6.}
\end{figure*}

\begin{figure*}[t]
\centering
\includegraphics[width=\textwidth]{../qual_figures/rhd_best/00143_qualitative.png}
\caption{Qualitative comparison on RHD evaluation example 00143.}
\end{figure*}

\begin{figure*}[t]
\centering
\includegraphics[width=\textwidth]{../qual_figures/rhd_best/00354_qualitative.png}
\caption{Qualitative comparison on RHD evaluation example 00354.}
\end{figure*}

\subsection{Differences in this implementation and DRHand}

\noindent\textbf{Point:} This repository is a DRHand-inspired adaptation rather than a strict reproduction.
\textbf{Evidence:} It keeps the same high-level dual-regression idea of a lightweight backbone, a heatmap branch, a coordinate branch, and a fusion rule, but the surrounding implementation and experiment design differ substantially from the original paper.
\textbf{Explanation:} The core method is recognizably DRHand-like, but the repo is not reproducing the original work line-for-line.
\textbf{Link:} The reported results should therefore be interpreted as evidence about this implementation rather than as a direct numerical replication of DRHand.

\medskip
\noindent\textbf{Point:} The dataset regime is materially different from the original DRHand paper.
\textbf{Evidence:} The original paper trains on YouTube2D Hands and GANerated Hands and evaluates on STB and RHD, whereas this repository focuses on RHD and a COCO-style hand dataset, and also includes a small manually annotated real-image pilot set.
\textbf{Explanation:} This changes both the training distribution and the types of generalisation questions the paper can answer.
\textbf{Link:} The empirical findings in this paper therefore apply to a different data setting from the original DRHand study.

\medskip
\noindent\textbf{Point:} This implementation broadens the training and evaluation pipeline beyond the original DRHand scope.
\textbf{Evidence:} The repo is implemented in PyTorch and adds configurable two-stage training, transfer-learning runs, a reduced 10-keypoint mode, shared-10 evaluation, MediaPipe baseline comparison, and deployment-oriented benchmarking on Jetson Orin Nano.
\textbf{Explanation:} These additions make the paper more focused on lightweight deployment, transfer, and practical HCI use than the original DRHand benchmark study.
\textbf{Link:} The present work should therefore be understood as an engineering extension of DRHand's core idea to a broader experimental setting.

\section{CONCLUSION AND FUTURE WORK}

\noindent\textbf{Point:} A reduced 10-keypoint representation is sufficient for the lightweight edge-HCI setting studied here.
\textbf{Evidence:} The reduced model remains close to the baseline in predictive quality while offering a runtime advantage.
\textbf{Explanation:} This shows that the full 21-keypoint output is not strictly necessary for the interaction-focused scope of this study.
\textbf{Link:} The first research question is answered positively.

\medskip
\noindent\textbf{Point:} Dual-branch fusion is justified, but only as a qualified improvement.
\textbf{Evidence:} Fusion gives the best overall balance across the reported results, but the diagnostic analysis shows that the current selection rule is not fully optimal.
\textbf{Explanation:} The two-branch design is useful, yet the present heuristic does not extract all of its potential benefit.
\textbf{Link:} The second research question is answered as a qualified yes.

\medskip
\noindent\textbf{Point:} Pseudo-labeled real-image data improves generalisation beyond synthetic-only training, especially when used for target-domain adaptation.
\textbf{Evidence:} The transfer results show that zero-shot cross-domain generalisation is weak, but fine-tuning on the target domain substantially improves performance.
\textbf{Explanation:} This means pseudo-labeled real-image supervision is helpful, but mainly as an adaptation mechanism rather than as a complete replacement for target-domain training.
\textbf{Link:} The third research question is also answered positively, with that limitation.

\medskip
\noindent\textbf{Point:} Future work should redesign the architecture specifically for reduced 10-keypoint prediction.
\textbf{Evidence:} The current results show that the reduced target remains effective, but the model still follows essentially the same DRHand-inspired design as the 21-keypoint version.
\textbf{Explanation:} A purpose-built reduced architecture may allow further reductions in model size and latency.
\textbf{Link:} The next step is therefore to test whether a smaller model can preserve the same practical accuracy in the 10-keypoint setting.

\medskip
\noindent\textbf{Point:} Real-image evaluation should be broadened beyond the current pilot set.
\textbf{Evidence:} The present study includes only a small manually annotated real-image dataset used for supplementary quantitative and qualitative comparison.
\textbf{Explanation:} A larger independently annotated real-image benchmark is needed to assess practical robustness more convincingly.
\textbf{Link:} Future work should therefore prioritise broader real-world evaluation.

\medskip
\noindent\textbf{Point:} Fusion remains one of the clearest opportunities for improvement.
\textbf{Evidence:} The current diagnostics show that the heuristic often fails to select the lower-error branch, and the original DRHand formulation also leaves room for stronger branch combination~\cite{dualreg}.
\textbf{Explanation:} A more principled or learned fusion strategy could improve both robustness and accuracy.
\textbf{Link:} Future work should therefore investigate improved fusion design.

\medskip
\noindent\textbf{Point:} Future work should introduce an explicit hand-detection stage to standardise the hand input before pose estimation.
\textbf{Evidence:} The current learned pipeline assumes a prepared hand region, while practical systems such as MediaPipe first localise the hand and then estimate landmarks from a normalized crop.
\textbf{Explanation:} A detection stage could produce more consistent hand scale, position, and framing, making the pose input distribution closer to the data seen during training and reducing reliance on aggressive spatial augmentation.
\textbf{Link:} Future work should therefore evaluate whether detector-based input normalisation improves robustness and simplifies the pose-learning stage.

\clearpage
\begin{thebibliography}{00}

\bibitem{vuletic}
T. Vuletic, A. Duffy, L. Hay, C. McTeague, and G. Campbell, ``Systematic literature review of hand gestures used in human-computer interaction interfaces,'' \emph{Int. J. Human-Computer Studies}, vol. 129, pp. 47--65, 2019.

\bibitem{wachs}
J. P. Wachs, M. Kolsch, H. Stern, and Y. Edan, ``Vision-based hand-gesture applications,'' \emph{Communications of the ACM}, vol. 54, no. 2, pp. 60--71, 2011.

\bibitem{li_vr}
Y. Li, J. Huang, F. Tian, H.-A. Wang, and G.-Z. Dai, ``Gesture interaction in virtual reality,'' \emph{Virtual Reality \& Intelligent Hardware}, vol. 1, no. 1, pp. 84--112, 2019.

\bibitem{dualreg}
D. Wei, S. An, X. Zhang, J. Tian, K. A. Tsintotas, A. Gasteratos, and H. Zhu, ``Dual regression for efficient hand pose estimation,'' in \emph{Proc. IEEE Int. Conf. Robotics and Automation}, 2022.

\bibitem{mediapipe}
F. Zhang, V. Bazarevsky, A. Vakunov, A. Tkachenka, G. Sung, C.-L. Chang, and M. Grundmann, ``MediaPipe hands: On-device real-time hand tracking,'' \emph{arXiv preprint arXiv:2006.10214}, 2020.

\bibitem{srhandnet}
Y. Wang, B. Zhang, and C. Peng, ``SRHandNet: Real-time 2D hand pose estimation with simultaneous region localization,'' \emph{IEEE Trans. Image Processing}, vol. 29, pp. 2977--2986, 2019.

\bibitem{wingloss}
Z.-H. Feng, J. Kittler, M. Awais, P. Huber, and X.-J. Wu, ``Wing loss for robust facial landmark localisation with convolutional neural networks,'' in \emph{Proc. IEEE Conf. Computer Vision and Pattern Recognition}, 2018, pp. 2235--2245.

\bibitem{fasthand}
S. An, X. Zhang, D. Wei, H. Zhu, J. Yang, and K. A. Tsintotas, ``FastHand: Fast monocular hand pose estimation on embedded systems,'' \emph{Journal of Systems Architecture}, vol. 122, p. 102361, 2022.

\bibitem{attentionlite}
N. Santavas, I. Kansizoglou, L. Bampis, E. Karakasis, and A. Gasteratos, ``Attention! A lightweight 2D hand pose estimation approach,'' \emph{IEEE Sensors Journal}, vol. 21, no. 10, pp. 11488--11501, 2021.

\bibitem{handobjectgcn}
Y. Yin, C. McCarthy, and D. Rezazadegan, ``Real-time lightweight 3D hand-object pose estimation using temporal graph convolution networks,'' in \emph{AI 2024: Advances in Artificial Intelligence}, LNCS 15443, 2024, pp. 243--255.

\bibitem{everydayegohands}
A. Prakash, R. Tu, M. Chang, and S. Gupta, ``3D hand pose estimation in everyday egocentric images,'' in \emph{Proc. ECCV}, 2024.

\bibitem{splitehand}
K. H. Yeh, T. W. Hsu, and M. Sun, ``SPLite Hand: Sparsity-aware lightweight 3D hand pose estimation,'' in \emph{Proc. AICCC 2025}, Tokyo, Japan, 2025, pp. 1--10.

\bibitem{stereojetson}
E. Williams-Linera and J. M. Ramirez-Cortes, ``Stereo vision system based on the NVIDIA Jetson Nano for real-time evaluation of yoga poses,'' in \emph{Proc. IEEE}, 2024, pp. 1--7.

\bibitem{edgepose}
L. Zhang et al., ``EdgePose: Real-time human pose estimation scheme for industrial scenes,'' \emph{IEEE Access}, vol. 12, pp. 156702--156716, 2024.

\bibitem{lightweightgcnhand}
Y. Zhang et al., ``A lightweight hand attitude estimation method based on GCN feature enhancement,'' \emph{Electronics}, vol. 13, no. 10, p. 1952, 2024.

\bibitem{handaugment}
Z. Zhang, S. Xie, M. Chen, and H. Zhu, ``HandAugment: A simple data augmentation method for depth-based 3D hand pose estimation,'' \emph{arXiv preprint arXiv:2001.00702}, 2020.

\bibitem{efficientpose}
I. B. Toft, J. P. B. Nielsen, N. K. Poulsen, and A. L. Laursen, ``EfficientPose: Scalable single-person pose estimation,'' \emph{arXiv preprint arXiv:2008.09062}, 2020.

\bibitem{blazepose}
V. Bazarevsky, I. Grishchenko, K. Raveendran, T. Zhu, F. Zhang, and M. Grundmann, ``BlazePose: On-device real-time body pose tracking,'' \emph{arXiv preprint arXiv:2006.10204}, 2020.

\end{thebibliography}

\end{document}
